% Source coverage: physical PDF pp. 101--102 (printed pp. 78--79).
% Contents: Problems 1--4; References 1--7, with Reference 6 continued on p. 102.
% Figures/tables/footnotes: none. Uncertainties: none after rendered-page review.

\chapterbackmatter{Problems}

\begin{enumerate}
  \item Consider a theory of real pseudoscalars $\phi(x)$ and complex Dirac fields
  $\psi(x)$, with masses $M$ and $m$ respectively, and interaction
  $g\bar\psi\gamma_5\psi\phi$. Evaluate the effective potential for $\phi={}$constant,
  $\psi=0$, to one-loop order.

  \item Derive general formulas for
  \begin{gather*}
    \frac{\delta^3W[J]}{\delta J_n(x)\delta J_m(y)\delta J_\ell(z)}
    \quad\text{and}\quad \\
    \frac{\delta^4W[J]}{\delta J_n(x)\delta J_m(y)\delta J_\ell(z)\delta J_k(w)}
  \end{gather*}
  in terms of the variational derivatives of $\Gamma[\phi]$ with respect to $\phi$.
  Show which Feynman diagrams correspond to each term in these formulas.

  \item Calculate the effective potential to one-loop order for the theory of a
  neutral scalar field $\phi$ with interaction Lagrangian density $g\phi^3/6$ in
  six spacetime dimensions.

  \item Suppose that the action $I[\phi]$ is invariant under a \emph{finite} matrix
  transformation $\phi_n(x)\longrightarrow\sum_m M_{nm}\phi_m(x)$. Under which
  transformation of the currents is $W[J]$ then invariant? Use this result to derive
  a symmetry property of $\Gamma[\phi]$.
\end{enumerate}

\chapterbackmatter{References}

\begin{enumerate}
  \item \phantomsection\label{ch16-ref-1}
  The effective action $\Gamma[\phi]$ was introduced by J.~Goldstone, A.~Salam,
  and S.~Weinberg, \textit{Phys. Rev.} \textbf{127}, 965 (1962), who defined it
  perturbatively, as a sum over one-particle-irreducible connected diagrams. The
  non-perturbative definition \eqref{eq:16.1.6} was first given by G.~Jona-Lasinio,
  \textit{Nuovo Cimento} \textbf{34}, 1790 (1964).

  \item \phantomsection\label{ch16-ref-2}
  S.~Coleman, \textit{Aspects of Symmetry} (Cambridge University Press, Cambridge,
  1985): pp.~135--6.

  \item \phantomsection\label{ch16-ref-3}
  S.~Coleman and E.~Weinberg, \textit{Phys. Rev.} D\textbf{7}, 1888 (1973).

  \item \phantomsection\label{ch16-ref-4}
  K.~Symanzik, \textit{Comm. Math. Phys.} \textbf{16}, 48 (1970); S.~Coleman,
  \textit{Aspects of Symmetry} (Cambridge University Press, Cambridge, 1985):
  pp.~139--42.

  \item \phantomsection\label{ch16-ref-5}
  K.~Symanzik, Ref.~4; J.~Iliopoulos, C.~Itzykson, and A.~Martin,
  \textit{Rev. Mod. Phys.} \textbf{47}, 165 (1975).

  \item \phantomsection\label{ch16-ref-6}
  Y.~Fujimoto, L.~O'Raifeartaigh, and G.~Parravicini, \textit{Nucl. Phys.}
  B\textbf{212}, 268 (1983); R.~W.~Haymaker and J.~Perez-Mercader,
  \textit{Phys. Rev.} D\textbf{27}, 1948 (1983); C.~M.~Bender and F.~Cooper,
  \textit{Nucl. Phys.} B\textbf{224}, 403 (1983); M.~Hindmarsh and D.~Johnston,
  \textit{J. Math. Phys.} A\textbf{19}, 141 (1986); V.~Branchina, P.~Castorina,
  and D.~Zappal\`a, \textit{Phys. Rev.} D\textbf{41}, 1948 (1990); K.~Cahill,
  \textit{Phys. Rev.} D\textbf{52}, 4704 (1995).

  \item \phantomsection\label{ch16-ref-7}
  A.~A.~Slavnov, \textit{Theor. Math. Phys.} \textbf{10}, 152 (1972)
  [English translation: \textit{Theor. and Math. Phys.} \textbf{10}, 99 (1972)];
  J.~C.~Taylor, \textit{Nucl. Phys.} B\textbf{33}, 436 (1971).
\end{enumerate}
